\section{引言}
一个先在菜谱应用里做菜、再去填购物清单的移动智能体，会在第 1 屏读到 ``vanilla extract''，并要在第 9 屏把它输进去——那时承载这段文字的屏幕早已离开显示区。我们把"所需文字曾出现在某个更早屏、且仅出现在那里"的步骤称为\textbf{延迟信息需求}；承载该文字的屏幕是\textbf{源屏}，其后经过的屏数是\textbf{距离}。这类步骤往往决定一个长跨应用任务能否完成，而用于研发移动智能体的基准几乎不含这类步骤。AITZ~\citep{aitz} 在 18,643 步中一例也没有，这与其任务平均仅 7.4 步相符；GUI-Odyssey~\citep{guiodyssey} 平均 15.4 步，抽样轨迹中有 31\% 至少含一例（\S\ref{sec:needs}）。

\paragraph{问题所在。} 有两种 prompt 式设计可以取回这段文字，而两者都随历史变长而退化。把每个更早屏的完整 OCR 文本都放进上下文，我们称之为\textbf{全历史}，它在第 15 步前出现的需求上定位正确率为 68\%，在更晚的需求上为 46\%。要点索引式的\textbf{屏幕索引}~\citep{readagent}——每屏写一行短摘要、由读取器据此挑一屏再回读原文——整体定位 41\%，而全历史为 54.8\%。我们用五种方式训练或重写其 writer，定位都停在原处；一个在采样行上取最优的 oracle 也只能达到全历史的水平，不能更高（\S\ref{sec:index}）。prompt 是两者共同的约束。据我们所知唯一训练过回看的工作 PAL-UI~\citep{palui} 需要 8.6K 条人工整理的逐步指令，因此在新语料上训练任一设计，都意味着要为"该回看哪一屏"付出标注代价。

\paragraph{关键观察。} 其实不必付。一个字符串检查定义了延迟信息需求，而这个检查同时回答了训练所需的问题。对第 $t$ 步的一个已验证需求，每个更早屏的 OCR 要么包含所输入的文字、要么不包含，于是挖掘该需求的验证器一次性给出了全部 $t$ 个候选的标签。我们把在这些标签上训练回看策略称为\textbf{验证器监督的回看}。这些标签的代价是每屏一次子串测试，且没有任何人查看过其中任何一条。

\paragraph{方法直觉。} 我们在验证器标签上训练两个实现。它们的差别在于历史如何抵达模型，并因此决定了同一批标签能产生多少个决策。\textbf{逐屏打分器}是一个带低秩适配器的 7B 模型，它每次读入任务、当前屏和一个更早屏，输出该更早屏含有所需文字的概率，定位即取各屏概率的最大者。\textbf{全历史读取器}是同一基座模型，在同样的标签上微调为在一次调用中从整段历史里直接输出屏号。随后由第二次调用即\textbf{读取阶段}产生要输入的文字，它或者只看选中的那一屏（\textbf{单屏读取}），或者看完整历史并把选中屏重复一遍加以标注（\textbf{带历史读取}）；\textbf{端到端}准确率评判最终输入的文字。

\paragraph{主要结果。} 起作用的是监督，不是架构。在按轨迹分组的五折交叉验证下，188 个已验证样例上，训练把全历史读取器从 54.8\% 提到 65.4\%（\pval{0.019}），把逐屏打分器从 50.0\% 提到 72.3\%（$p<0.001$）。两个训练变体之间定位相差 6.9 个点、端到端相差 2.1 个点，配对检验对二者均不可区分（\pval{0.086} 与 \pval{0.45}）。打分器的增益集中在 8--12 屏的距离区间，那里它定位 78.2\%，而训练后的全历史读取器为 56.4\%（\pval{0.003}）。我们在冻结全部设计决策之后才挖掘第四批样本，并只评测一次：打分器定位 79.6\%，训练后的全历史读取器 76.3\%（\pval{0.63}），prompt 全历史 59.1\%（\pval{0.003}）；带历史读取把端到端准确率从 23.7\% 提到打分器的 49.5\% 与全历史读取器的 51.6\%（两者均 $p<0.001$）。在 13 屏及以上的距离上，两个训练变体都跌到 prompt 全历史之下。

\paragraph{贡献。}（1）\textbf{问题与标签。} 延迟信息需求的定义、一个同时完成挖掘与标注的自动验证器、该需求的普遍程度，以及任何相关研究都必须施加的两个过滤条件（\S\ref{sec:needs}）。（2）\textbf{方法。} 验证器监督的回看，以及共享同一训练目标的两个实现（\S\ref{sec:method}）。（3）\textbf{发现。} 监督才是杠杆：它同时抬升两种架构，而逐屏分解带来的额外增益被限定在中等距离区间（\S\ref{sec:main}、\S\ref{sec:gap}）。（4）\textbf{验证与边界。} 一批外部样本确认了相对 prompt 的优势，消融表明该设计对其 prompt 与损失设置不敏感，而 13 屏及以上的距离仍未解决（\S\ref{sec:ext}、\S\ref{sec:limits}）。
